%% campaign.tex — sample campaign book for apotheosis-book.cls
%% Compile with:  latexmk   (uses XeLaTeX via latexmkrc)
%%            or: xelatex campaign.tex  (twice, for the contents)
%%
%% Themes: dark (default) · vermis · light · print
\documentclass{apotheosis-book} % or [vermis], [light], [print], [a4paper]

%% ================ TITLE ================ %%
\title{Apotheosis}
\subtitle{A Campaign of Lyrium, Blood and Ascension}
\author{The Charred Eye}
\version{0.1 (template preview)}
\coverimage{art/cover.jpg}
\credits{
  Template built on the Dragonbane Supplement class by Sibling Dex and the
  D\&D 5e LaTeX Template by the rpgtex contributors.

  Interior art from the Apotheosis campaign notes.

  Layout inspired by the official fifth edition books and by
  Plastiboo's \emph{Vermis}.
}

\begin{document}
\maketitle
\makecredits
\tableofcontents

%% ================================================================ %%
\partimage{art/lich.jpg}
\part{The Player's Codex}

\chapterart{art/red-wastes.png}
\chapter{Changes to the Base Game}

\begin{quote}
  As you will soon discover, this world is ripe with a feast of options and
  varieties of character flavour heard before only on the lips of a madman.
  \quoteby{The Charred Eye}
\end{quote}

\begin{segment}
\dropcap{T}{hrough the powers} that be, this campaign has additional
mechanics for the adventurers involved. From elven princesses of argent
starlight to clowns of eldritch fear, this guide to \textbf{Apotheosis}
outlines your unique powers and afflictions.

\section{Homebrew Rules}
This section outlines universal changes to the rules, from combat to
travel and even roleplay.

\subsection{Travel}
Travel is broken into three tiers of distance and cost. Depending on how
fast you go and the route you take, encounters can be avoided and the
resources you find can grow or shrink.

\begin{rolltable}[Travel Tiers]{l c L}
  \colhead{Distance} & \colhead{Encounters} & \colhead{Rest Afterwards} \\
  Nearby   & 1+ & Short rest \emph{recommended} \\
  Far      & 2+ & Short rest \emph{required} \\
  Very Far & 3+ & Long rest \emph{required} \\
\end{rolltable}

\subsection{Class Divergence}
Some classes may use a \textbf{variant spellcasting ability} that fits the
character better.

\subsubsection{Paladin}
The spellcasting modifier for the paladin is \textbf{Wisdom} instead of
\textbf{Charisma}.

\subsubsection{Warlock}
The base warlock uses \textbf{Charisma}, though you may ask to use
\textbf{Intelligence} instead.

\subsubsection{Sorcerer}
The base sorcerer uses \textbf{Charisma}; if you never multiclass out of
sorcerer, you may ask to use \textbf{Constitution}.

\begin{readaloud}
  The bell of the drowned chapel rings once, though no rope hangs from it.
  Somewhere beneath the black water, something that was once a knight
  begins to climb.
\end{readaloud}

\subsection{Critical Rolls}
When a \textbf{natural 20} or \textbf{natural 1} is rolled, the critical
success or failure gains an added effect chosen by the Dungeon Master.

\columnart[height=2.6in, focus=.3]{art/dreaded-lady.jpg}

\begin{sidebar}{Two Halves of a Whole}
  An organised whole that is perceived as more than the sum of its parts.
  Wielders of blade \emph{and} sorcery walk a narrow road, and the road
  remembers every step.
\end{sidebar}

\subsection{Hybrid Source Material}
This campaign uses the \emph{2014} and \emph{2024 Player's Handbooks}
together, to widen the roleplaying potential of the adventurers and the
world built around them.

\paragraph{Species \& Backgrounds.} See the relevant chapter for the
custom and adapted species available.
\end{segment}

\begin{variantrule}{Heroic Death \& Act of Valour}
\begin{segment}
  \paragraph{Act of Valour.} Once per session, when an ally within 30 feet
  would drop to 0 hit points, you can take the blow instead.

  \paragraph{Heroic Death.} A character who dies furthering the party's
  goal may make one final action of their choice before falling, and the
  table remembers it.
\end{segment}
\end{variantrule}

%% ================================================================ %%
\chapterart[focus=.6]{art/blue-storm.png}
\chapter{Afflictions of Lyrium}

\begin{segment}
\dropcap{L}{yrium sings.} Blue lyrium sings softly, of memory and of
power; red lyrium screams. Every character who handles it long enough
learns the tune.

\spellheader{Caeruleolyriasis}
  {Special affliction (blue lyrium withdrawal)}
  {1 hour (escalates)}{Self}{None}{Until cured}
When you are separated from a steady supply of \textbf{blue lyrium}, or
try to stop using it, you suffer \textbf{Caeruleolyriasis}. While
afflicted you feel fatigue, forgetfulness, a gnawing thirst and cold
hands, and the hunger only deepens.

\begin{DndTable}[header=Withdrawal Side Effects]{c L}
  \colhead{d6} & \colhead{Effect} \\
  1 & \textbf{Dermal Luminescence.} You shed dim light 5 ft.\ and have
      disadvantage on Dexterity (Stealth). \\
  2 & \textbf{Arcano-Dyslexia.} For 1 hour, succeed on a DC 12
      Intelligence check to read normally. \\
  3 & \textbf{Perceptual Desync.} Disadvantage on initiative rolls. \\
  4 & \textbf{Fine-Motor Tremor.} Disadvantage on Dexterity checks made
      with tools. \\
  5 & \textbf{Scar Irradiation.} You can't gain temporary hit points. \\
  6 & \textbf{Resistance.} You resist further withdrawal effects. \\
\end{DndTable}

\subsection{Signs of Exposure}
\begin{itemize}
  \bolditem{Voices:} A faint Calling, like a choir heard through stone.
  \coloritem{Crystals:} Red veins that surface under the skin.
  \coloritem[Frost]{Cold Hands:} The first sign of blue withdrawal.
  \secretitem{The Truth:} Red lyrium is alive, and it is listening.
\end{itemize}

\end{segment}

%% ================================================================ %%
\partcolor{Lyrium}
\partimage{art/rotten-one.jpg}
\part{The Bestiary}

\chapterart[focus=.25]{art/dhampir.jpg}
\chapter{Things That Wait Below}

\begin{segment}
\dropcap{B}{eneath the Tower} of Elves, the lyrium veins grow thick enough
to walk on. What lives there was not born; it was \emph{grown}.

\begin{monster}{Behemoth}
\monstertype{Large monstrosity, unaligned}
\basics[
  armorclass = 10 (lyrium hide),
  hitpoints  = \rpgdice{16d10+32},
  speed      = 35 ft.
]
\monsterrule
\stats[
  STR = \stat{26}, DEX = \stat{10}, CON = \stat{26},
  INT = \stat{8},  WIS = \stat{12}, CHA = \stat{8}
]
\monsterrule
\details[
  skills    = Athletics +10,
  senses    = passive Perception 11,
  languages = understands Common but can't speak,
  challenge = 10
]
\monsterrule
\begin{monsteraction}[Essence of Red Lyrium]
  The behemoth has advantage on saving throws against spells and other
  magical effects.
\end{monsteraction}
\begin{monsteraction}[Mage Slayer]
  When the behemoth damages a spellcaster, the target must succeed on a
  DC 16 Constitution saving throw or lose its lowest spell slot.
\end{monsteraction}

\monstersection{Actions}
\begin{monsteraction}[Bash]
  \emph{Melee Weapon Attack:} +10 to hit, reach 10 ft., one target.
  \emph{Hit:} \rpgdice{2d6+8} bludgeoning damage.
\end{monsteraction}
\begin{monsteraction}[Wall of Lyrium]
  Jagged red lyrium bursts from the ground at up to four points within
  30 ft. Each creature within 5 ft.\ of a point must succeed on a DC 13
  Dexterity saving throw or take \rpgdice{2d6} piercing damage.
\end{monsteraction}

\monstersection{Bonus Actions}
\begin{monsteraction}[Quake]
  The behemoth spends 2 \textbf{Shards}. Each creature of its choice within
  10 ft.\ must succeed on a DC 16 Strength saving throw or fall prone.
\end{monsteraction}
\end{monster}

\section{The Drowned Chapel}
\area{The Nave}
Black water fills the nave to the knee. Pews float, overturned, and the
altar glows a faint and hungry red.

\subarea{The Font}
The font is dry. Something has scratched \emph{``she is still singing''}
into the stone, over and over.

\area{The Bell Tower}
A narrow stair spirals up into the dark. The bell has no rope.

\columnart[fade=all, height=2.4in]{art/black-knight.jpg}

\begin{note}{Shards}
  Lyrium-touched creatures spend \textbf{Shards} to fuel their abilities.
  Track them with tokens on the table where players can see them dwindle.
\end{note}
\end{segment}

\fullpageart{art/black-knight.jpg}[The Black Knight does not sleep. It waits.]

%% ================================================================ %%
\setaccent{Ember}
\chapter*{Template Reference}
\addcontentsline{toc}{chapter}{Template Reference}
\pagebottomart{art/red-wastes.png}

\begin{segment}
\section{Art That Blends In}
\begin{description}
  \item[\textbackslash coverimage\{file\}] Full-bleed cover behind \texttt{\textbackslash maketitle}.
  \item[\textbackslash partimage\{file\}] Full-bleed art for the next part page.
  \item[\textbackslash chapterart{[keys]}\{file\}] Art across the top of the next chapter opener.
  \item[\textbackslash columnart{[keys]}\{file\}] Soft-edged art inside a column.
  \item[\textbackslash pagetopart / \textbackslash pagebottomart] Full-bleed bands on the current page.
  \item[\textbackslash fullpageart\{file\}{[caption]}] A page that is only art.
  \item[\textbackslash fadedart{[keys]}\{file\}] The raw box behind all of them.
\end{description}
Keys: \texttt{fade} (soft, all, bottom, top, sides, left, right, none),
\texttt{top}, \texttt{bottom}, \texttt{left}, \texttt{right} (0--0.5),
\texttt{width}, \texttt{height}, \texttt{focus} (0--1) and \texttt{opacity}.

\section{Boxes}
\texttt{readaloud}, \texttt{sidebar}, \texttt{variantrule}, \texttt{note},
\texttt{monster}, \texttt{rolltable} and \texttt{tablebox}. The Dragonbane
names (\texttt{dragonbox}, \texttt{demonbox}, \texttt{emptybox}) and the D\&D
names (\texttt{DndReadAloud}, \texttt{DndSidebar}, \texttt{DndTable},
\texttt{monsterbox}, \texttt{\textbackslash DndSpellHeader}) still work, so
existing notes move over with few edits.

\section{Colour}
\texttt{\textbackslash partcolor\{Name\}} or \texttt{\textbackslash setaccent\{Name\}}
with Ember, Coral, Ochre, Blood, Bile, Mustard, Frost, Amethyst or Lyrium.
\end{segment}

\end{document}
